% AUTO-GENERATED from scripts_pa/REPORTING_CHECKLIST_V2.md (v2 draft, 2026-10-22)
% do not edit: this file is generated
% v1 (scripts_pa/REPORTING_CHECKLIST.md) is unchanged and still shipped as Checklist S1

\paragraph{A. Graph construction visibility}
\begin{enumerate}
\setcounter{enumi}{0}
\item[{\textbf{1.}}] \textbf{[must]} Are the outer test subjects present as nodes in the graph used during training (parameter learning)? (\texttt{R0-T} vs \texttt{R1-C} distinction)
\item[{\textbf{2.}}] \textbf{[must]} Are outer test labels excluded from every fitting step (feature selection, normalization statistics, early stopping, threshold selection, hyperparameter search)?
\item[{\textbf{3.}}] \textbf{[if]} If test subjects are present at training time without labels, is this stated explicitly as \emph{unlabeled cohort visibility} rather than \emph{label leakage}?
\item[{\textbf{4.}}] \textbf{[must]} Are graph edges (or edge-candidate sets) constructed from the same subject set in training and inference? If not, name both sets.
\item[{\textbf{5.}}] \textbf{[must]} Are phenotypic variables of test subjects (e.g. acquisition site, sex, age) used in edge construction? Report yes/no per variable.
\end{enumerate}

\paragraph{B. Query context at inference}
\begin{enumerate}
\setcounter{enumi}{5}
\item[{\textbf{6.}}] \textbf{[must]} Is a query subject's prediction computed with other unlabeled test subjects co-present in the inference graph (cohort-visible) or in isolation (single-query)?
\item[{\textbf{7.}}] \textbf{[if]} If cohort-visible, report the context size convention and whether it is the full remaining test cohort, a fixed batch size \texttt{k}, or sampled.
\item[{\textbf{8.}}] \textbf{[if]} If two inference contexts are compared (for example the full remaining test cohort versus a single query), is the checkpoint identical between them, and is its hash reported?
\item[{\textbf{9.}}] \textbf{[must]} Is the inference graph permutation invariant with respect to the query's node index? If not, state the fixed node-order convention and report the effect of the alternative convention.
\item[{\textbf{10.}}] \textbf{[must]} Is the decision threshold selected without using query labels?
\end{enumerate}

\paragraph{C. Site familiarity}
\begin{enumerate}
\setcounter{enumi}{10}
\item[{\textbf{11.}}] \textbf{[must]} Report the acquisition-site composition of the inference context (same-site vs cross-site vs mixed) and, when possible, the size-matched same-site vs cross-site contrast.
\item[{\textbf{12.}}] \textbf{[if]} If leave-one-site-out (or unseen-site) results are reported, state that neither features nor labels of the held-out site were used at any fitting step.
\item[{\textbf{13.}}] \textbf{[must]} Do not describe a model that saw held-out-site features at training time as "unseen-site".
\end{enumerate}

\paragraph{D. Normalization and preprocessing scope}
\begin{enumerate}
\setcounter{enumi}{13}
\item[{\textbf{14.}}] \textbf{[must]} Report the scope (fit set) of every normalization or harmonization step: full cohort before splitting, fold-local, or F-only.
\item[{\textbf{15.}}] \textbf{[if]} If any normalization statistic depends on the current inference graph (cohort-dependent normalization), report it as an information channel and, where possible, provide a frozen-statistics sensitivity run.
\end{enumerate}

\paragraph{E. Estimation and inference}
\begin{enumerate}
\setcounter{enumi}{15}
\item[{\textbf{16.}}] \textbf{[must]} State the primary estimand, its unit of analysis (subject, site, partition) and what would make the effect scientifically meaningful. If individual-level prediction shifts are part of the claim, report an appropriate shift measure and its aggregation rule. If a smallest effect size of interest is invoked, state its basis (substantive, prior evidence, or decision context) and the date it was fixed; a threshold derived from the current estimate is reported only as an exploratory sensitivity analysis.
\item[{\textbf{17.}}] \textbf{[must]} Report paired subject-level inference (bootstrap or permutation) using the same resample for both regimes.
\item[{\textbf{18.}}] \textbf{[must]} Cluster resampling by acquisition site when subjects are nested within sites (two-stage site$\to$subject), and compare with subject-level bootstrap as a sensitivity analysis.
\item[{\textbf{19.}}] \textbf{[must]} State the scope of simultaneous inference. If the claim requires a family-wise or simultaneous statement, pre-specify the family and the correction (e.g. Holm) and report raw and adjusted quantities. If no simultaneous claim is made, say so explicitly and report intervals as marginal, per-partition intervals.
\item[{\textbf{20.}}] \textbf{[if]} If a model-by-regime interaction is claimed, estimate the difference-in-differences directly and report uncertainty that preserves the pairing (same resample for both regimes, and site clustering where subjects are nested in sites). If a hypothesis test is used, name the test, the null it samples under, and the multiplicity handling. A point estimate without any uncertainty statement is not an interaction claim; a point estimate with a paired interval suited to the estimand and to its dependence structure is.
\item[{\textbf{21.}}] \textbf{[must]} Separate split-seed variance from subject-level sampling variance; never treat seeds or context draws as independent subjects.
\end{enumerate}

\paragraph{F. Reproducibility}
\begin{enumerate}
\setcounter{enumi}{21}
\item[{\textbf{22.}}] \textbf{[must]} Provide the cohort manifest, split registry (with hash), model commit, environment lock, and per-cell checkpoint hashes.
\item[{\textbf{23.}}] \textbf{[must]} Provide per-subject out-of-fold predictions in a documented schema (subject, site, label, split seed, fold, train seed, model, regime, prediction), so that all reported numbers can be regenerated mechanically.
\item[{\textbf{24.}}] \textbf{[must]} State the node-order convention, the context sampling rule, and the threshold rule in the methods section, not only in code.
\end{enumerate}

\paragraph{G. Claims}
\begin{enumerate}
\setcounter{enumi}{24}
\item[{\textbf{25.}}] \textbf{[must]} State the deployment assumption that the reported regime corresponds to. Cohort-visible inference answers a \emph{cohort-conditioned} question; single-query inference answers a \emph{prospective} question. Do not present either as the only correct protocol.
\item[{\textbf{26.}}] \textbf{[must]} Do not state that a model's absolute accuracy is \emph{explained by} site composition unless a site-specific intervention supports that claim.
\item[{\textbf{27.}}] \textbf{[must]} Report the number of excluded sites/subjects and the exclusion rule.
\end{enumerate}

\subsection*{Change record relative to v1}
Items 8, 16, 19 and 20 are revised. Item 8 becomes conditional ([if]) and is recorded as not applicable here, because no query-isolated inference contrast is reported; items 16, 19 and 20 are rewritten, and items 1--7, 9--15, 17--18 and 21--27 are unchanged in wording and numbering. The reason for the revisions is narrow: in v1, items 16, 19 and 20 treated one statistical route (hypothesis testing with a pre-specified SESOI and corrected $p$-values) as a universal requirement, so an estimation-based study could not satisfy them whatever its evidence; item 8 assumed a two-context comparison that a single-context study cannot answer. v2 keeps both statistical routes admissible and asks that the route be declared, and makes item 8 apply only when such a comparison is made. The change is not a relaxation made to let this study pass: the same wording applies to any study, and under v1 this study would still be recorded as not meeting items 16, 19 and 20 (that record is kept in Checklist S1).

\subsection*{Applicability and evidence pointers, item by item}
Status is stated against v2. ``Met (v2)'' marks the three items whose assessment changes because the requirement changed; no other item changes status.

\begin{description}
\item[1.] Met --- \S2.3; R0-T/R1-C defined
\item[2.] Met --- \S2.2 (fit-fold selection); OOF registry
\item[3.] Met --- \S2.3 terminology
\item[4.] Met --- \S2.2; split registry (\S2.5)
\item[5.] Met --- \S2.2 (phenotype only in edge weights)
\item[6.] Met --- \S2.3; both regimes cohort-visible
\item[7.] Partial --- inference context is the remaining test cohort (\S2.3); the size convention is not stated in the main text
\item[8.] Not applicable --- no query-isolated inference contrast is reported; the inference axis is excluded (\S2.3)
\item[9.] Partial --- fixed node order with the query last (\S2.2); no permutation comparison is reported
\item[10.] Met --- \S2.2 threshold rule (Youden's $J$ on the validation partition only)
\item[11.] Partial --- site composition of each split is in the registry; the inference-context composition is not tabulated
\item[12.] Not applicable --- no held-out-site training regime reported
\item[13.] Met --- no unseen-site claim
\item[14.] Met --- \S2.2
\item[15.] Not applicable --- no cohort-dependent normalization
\item[16.] Met (v2) --- \S2.4 estimand; no per-subject shift claimed
\item[17.] Met --- \S2.6 paired bootstrap
\item[18.] Met --- \S2.6 site-cluster + subject bootstrap
\item[19.] Met (v2) --- \S2.6 marginal intervals declared
\item[20.] Met (v2) --- \S3.1--3.2 Delta-Delta AUC with paired intervals
\item[21.] Met --- \S2.6 seeds vs subjects
\item[22.] Partial --- 1024/2024 yes; split 666 no
\item[23.] Partial --- 1024/2024 yes; split 666 no
\item[24.] Partial --- 1024/2024 yes; split 666 not documented
\item[25.] Met --- \S2.3 deployment assumption stated
\item[26.] Met --- no site-composition explanation claimed
\item[27.] Partial --- cohort size and site count are stated (\S2.1); the exclusion rule is only in the released repository
\end{description}
